\chapter{Chuỗi Markov thời gian liên tục\\{\normalsize\textit{Continuous-Time Markov Chains}}}
\label{ch:continuous}

% ============================================================================
\section{Quá trình Poisson (The Poisson Process)}
\label{sec:poisson}

\textbf{Câu chuyện:} Hãy tưởng tượng bạn ngồi ở quán café quan sát xe buýt đến trạm.
Trung bình cứ 10 phút có một chuyến, nhưng chính xác khi nào thì xe đến --- bạn không biết trước.
Có khi 5 phút đã có xe, có khi phải đợi 15 phút. Đây chính là bản chất của \term{quá trình Poisson}:
sự kiện xảy ra ``ngẫu nhiên'' theo thời gian, với một tốc độ trung bình cố định.

Trong chương này ta bắt đầu nghiên cứu các \term{quá trình thời gian liên tục} (\en{continuous-time processes}),
là các họ $(X_t)_{t\in\R_+}$ gồm các biến ngẫu nhiên được đánh chỉ số bởi $\R_+$.
Khác với chuỗi Markov rời rạc (thời gian nhảy tại $n = 0, 1, 2, \ldots$),
ở đây thời gian chạy liên tục và quá trình có thể ``nhảy'' vào bất kỳ lúc nào.

\begin{tomtat}
\textbf{Ý tưởng chính:} Quá trình Poisson đếm số sự kiện xảy ra theo thời gian.
Nó là ``viên gạch nền tảng'' cho mọi chuỗi Markov liên tục,
giống như bước đi ngẫu nhiên là nền tảng cho chuỗi Markov rời rạc.
\end{tomtat}

\begin{dinhnghia}[(Quá trình Poisson chuẩn -- \en{Standard Poisson process})]
\term{Quá trình Poisson chuẩn} (\en{standard Poisson process}) là quá trình ngẫu nhiên $(N_t)_{t\in\R_+}$ chỉ có bước nhảy $+1$, và đường đi (\en{paths}) không đổi giữa các bước nhảy. Tại thời điểm $t$:
\[
  N_t = \sum_{k=1}^{\infty} \ind_{(T_k, \infty)}(t), \qquad t \in \R_+,
\]
trong đó $(T_k)_{k\geq 1}$ là dãy tăng các thời điểm nhảy với $\lim_{k\to\infty} T_k = +\infty$.

$(N_t)_{t\in\R_+}$ thỏa mãn:
\begin{enumerate}[nosep]
  \item \textbf{Gia số độc lập} (\en{independence of increments}): với mọi $0 \leq t_0 < t_1 < \cdots < t_n$, các biến $N_{t_1} - N_{t_0}, \ldots, N_{t_n} - N_{t_{n-1}}$ là độc lập.
  \item \textbf{Gia số dừng} (\en{stationarity of increments}): $N_{t+h} - N_{s+h}$ có cùng phân phối với $N_t - N_s$ với mọi $h > 0$ và $0 \leq s \leq t$.
\end{enumerate}
\end{dinhnghia}

\textbf{Ví dụ đời thường:}
\begin{itemize}[nosep]
  \item Số cuộc gọi đến tổng đài trong khoảng thời gian $[0, t]$.
  \item Số khách hàng vào cửa hàng từ 8h sáng đến giờ $t$.
  \item Số hạt phóng xạ phát ra bởi một mẫu chất trong thời gian $t$.
\end{itemize}

\begin{dinhly}[(Định lý 9 -- Phân phối Poisson)]
Giả sử quá trình đếm $(N_t)_{t\in\R_+}$ thỏa mãn các điều kiện (1) và (2). Khi đó:
\[
  \PP(N_t - N_s = k) = e^{-\lambda(t-s)} \frac{(\lambda(t-s))^k}{k!}, \qquad k \in \N, \quad 0 \leq s \leq t,
\]
với hằng số $\lambda > 0$.
\end{dinhly}

Tham số $\lambda > 0$ được gọi là \term{cường độ} (\en{intensity}) của quá trình:
\begin{equation}\label{eq:intensity}
  \lambda := \lim_{h \to 0} \frac{1}{h} \PP(N_h = 1).
\end{equation}

Đặc biệt, $N_t$ có \term{phân phối Poisson} (\en{Poisson distribution}) với tham số $\lambda t$:
\[
  \PP(N_t = k) = \frac{(\lambda t)^k}{k!}\, e^{-\lambda t}, \qquad t > 0.
\]

\textbf{Xác suất chuyển vi mô} (\en{infinitesimal transition probabilities}):
\begin{align}
  \PP(N_{t+h} - N_t = 1) &\simeq \lambda h, \qquad h \to 0, \label{eq:poisson-inf-1}\\
  \PP(N_{t+h} - N_t = 0) &\simeq 1 - \lambda h, \qquad h \to 0. \label{eq:poisson-inf-0}
\end{align}

Điều này có nghĩa: trong khoảng thời gian ``ngắn'' $[t, t+h]$, gia số $N_{t+h} - N_t$ xử sự như biến ngẫu nhiên Bernoulli với tham số $\lambda h$.

\begin{proposition}[(Phân phối thời gian nhảy)]
Luật của $T_n$ có mật độ $t \mapsto \lambda^n e^{-\lambda t} \frac{t^{n-1}}{(n-1)!}$ trên $\R_+$, $n \geq 1$.
\end{proposition}

Thời gian $\tau_k := T_{k+1} - T_k$ giữa các bước nhảy tạo thành dãy các biến ngẫu nhiên \term{độc lập cùng phân phối} (\en{i.i.d.}) có phân phối mũ với tham số $\lambda > 0$:
\[
  \PP(\tau_0 > t_0, \ldots, \tau_n > t_n) = e^{-\lambda(t_0 + \cdots + t_n)} = \prod_{k=0}^n e^{-\lambda t_k},
\]
và $\E[\tau_k] = 1/\lambda$.

\begin{example}[Xe buýt đến trạm]
Giả sử xe buýt đến trạm theo quá trình Poisson với cường độ $\lambda = 6$ chuyến/giờ (trung bình 10 phút/chuyến).

\textbf{Câu hỏi 1:} Xác suất không có xe nào trong 15 phút?
\[
  \PP(N_{0.25} = 0) = e^{-6 \times 0.25} = e^{-1.5} \approx 0.223.
\]

\textbf{Câu hỏi 2:} Xác suất có đúng 2 xe trong 20 phút?
\[
  \PP(N_{1/3} = 2) = e^{-2} \frac{2^2}{2!} = 2e^{-2} \approx 0.271.
\]

\textbf{Câu hỏi 3:} Thời gian đợi trung bình cho xe tiếp theo?
\[
  \E[\tau] = \frac{1}{\lambda} = \frac{1}{6} \text{ giờ} = 10 \text{ phút}.
\]
\end{example}

\begin{tomtat}
Quá trình Poisson: gia số độc lập và dừng, $N_t \sim \text{Poisson}(\lambda t)$.
Cường độ $\lambda$ cao hơn $\Rightarrow$ nhảy thường xuyên hơn $\Rightarrow$ thời gian ở mỗi trạng thái ngắn hơn ($\E[\tau_k] = 1/\lambda$).

Tính chất ``không nhớ'' (\en{memoryless}): dù bạn đã đợi 5 phút rồi, thời gian đợi thêm vẫn có phân phối mũ --- quá khứ không ảnh hưởng tương lai.
\end{tomtat}

% ============================================================================
\section{Quá trình sinh-tử (Birth and Death Processes)}
\label{sec:birth-death}

\textbf{Câu chuyện:} Hãy tưởng tượng một quán café. Khách đến (``sinh'') và khách đi (``tử'').
Nếu quán đang có $i$ khách:
\begin{itemize}[nosep]
  \item Khách mới đến với tốc độ $\lambda_i$ (số khách tăng lên 1).
  \item Một khách hiện tại rời đi với tốc độ $\mu_i$ (số khách giảm đi 1).
\end{itemize}
Số khách trong quán thay đổi liên tục theo thời gian --- đây là \term{quá trình sinh-tử}.

\begin{dinhnghia}[(Quá trình sinh thuần túy -- \en{Pure birth process})]
Chuỗi Markov thời gian liên tục $(X_t)_{t\in\R_+}$ thỏa mãn:
\[
  \PP(X_{t+h} - X_t = 1 \mid X_t = i) \simeq \lambda_i h, \qquad h \to 0,\; i \in S,
\]
\[
  \PP(X_{t+h} - X_t = 0 \mid X_t = i) \simeq 1 - \lambda_i h, \qquad h \to 0,\; i \in S,
\]
được gọi là \term{quá trình sinh thuần túy} (\en{pure birth process}) với tốc độ sinh $\lambda_i \geq 0$.
Quá trình Poisson là trường hợp đặc biệt với $\lambda_n = \lambda$ không phụ thuộc trạng thái.
\end{dinhnghia}

\begin{dinhnghia}[(Quá trình tử thuần túy -- \en{Pure death process})]
Tương tự, $(X_t)_{t\in\R_+}$ với:
\[
  \PP(X_{t+h} - X_t = -1 \mid X_t = i) \simeq \mu_i h, \quad
  \PP(X_{t+h} - X_t = 0 \mid X_t = i) \simeq 1 - \mu_i h,
\]
là \term{quá trình tử thuần túy} (\en{pure death process}) với tốc độ tử $\mu_i \geq 0$.
\end{dinhnghia}

\begin{dinhnghia}[(Quá trình sinh-tử -- \en{Birth and death process})]
Chuỗi Markov thời gian liên tục $(X_t)_{t\in\R_+}$ với:
\begin{align*}
  \PP(X_{t+h} - X_t = 1 \mid X_t = i)  &\simeq \lambda_i h, \\
  \PP(X_{t+h} - X_t = -1 \mid X_t = i) &\simeq \mu_i h, \\
  \PP(X_{t+h} - X_t = 0 \mid X_t = i)  &\simeq 1 - (\lambda_i + \mu_i)h,
\end{align*}
khi $h \to 0$, được gọi là \term{quá trình sinh-tử} (\en{birth and death process}) với tốc độ sinh $\lambda_i \geq 0$ và tốc độ tử $\mu_i \geq 0$, $i \in S$.
\end{dinhnghia}

Thời gian $\tau_i$ ở trạng thái $i$ trước khi chuyển sang trạng thái khác là biến ngẫu nhiên có phân phối mũ:
\[
  \tau_i = \min(\tau_{i,i+1},\; \tau_{i,i-1}),
\]
với tham số $\lambda_i + \mu_i$, và $\E[\tau_i] = \frac{1}{\lambda_i + \mu_i}$.

\begin{example}[Quán café có sức chứa $N = 4$]
Xét quán café có tối đa $4$ chỗ ngồi. Khách đến với tốc độ $\lambda = 2$ người/giờ,
khách rời đi với tốc độ $\mu = 3$ người/giờ (mỗi khách ngồi trung bình 20 phút).

Khi quán có $i$ khách ($0 \leq i \leq 4$):
\begin{itemize}[nosep]
  \item Tốc độ sinh: $\lambda_i = \lambda = 2$ (nếu $i < 4$), $\lambda_4 = 0$ (quán đầy).
  \item Tốc độ tử: $\mu_i = i \cdot \mu = 3i$ (mỗi khách rời đi độc lập với tốc độ $\mu$).
\end{itemize}

Thời gian trung bình ở trạng thái $i = 2$:
\[
  \E[\tau_2] = \frac{1}{\lambda_2 + \mu_2} = \frac{1}{2 + 6} = \frac{1}{8} \text{ giờ} = 7.5 \text{ phút}.
\]
\end{example}

% ============================================================================
\section{Chuỗi Markov thời gian liên tục (Continuous-Time Markov Chains)}
\label{sec:ctmc-definition}

\textbf{Câu chuyện:} Giống như con kiến trên bàn cờ (Chương~\ref{ch:discrete}) chỉ quan tâm ô hiện tại,
chuỗi Markov liên tục cũng ``không nhớ'' --- nhưng bây giờ con kiến di chuyển trong thời gian liên tục.
Nó ở một ô trong thời gian ngẫu nhiên (phân phối mũ), rồi nhảy sang ô khác.

\begin{dinhnghia}[(Tính chất Markov liên tục -- \en{Continuous-time Markov property})]
Quá trình ngẫu nhiên $(X_t)_{t\in\R_+}$ nhận giá trị trong $\Z$ có \term{tính chất Markov} nếu với mọi
$0 < s_1 < \cdots < s_{n-1} < s < t$:
\[
  \PP(X_t = j \mid X_s = i_n,\; X_{s_{n-1}} = i_{n-1},\; \ldots,\; X_{s_1} = i_0) = \PP(X_t = j \mid X_s = i_n).
\]
\end{dinhnghia}

\textbf{Ví dụ:} Quá trình Poisson là chuỗi Markov thời gian liên tục vì nó có gia số độc lập.
Tổng quát hơn, mọi quá trình có gia số độc lập đều là chuỗi Markov.

\subsection{Nửa nhóm chuyển (Transition Semigroup)}

Xác suất chuyển phụ thuộc thời gian:
\[
  P_{i,j}(t) := \PP(X_{t+s} = j \mid X_s = i), \qquad i,j \in \Z,\; t \in \R_+,
\]
trong đó ta giả sử quá trình là \term{đồng nhất theo thời gian} (\en{time homogeneous}): $\PP(X_{t+s} = j \mid X_s = i)$ không phụ thuộc $s$. Lưu ý $P(0) = I_d$.

Ma trận $P(t) = [P_{i,j}(t)]_{i,j\in S}$ được gọi là \term{nửa nhóm chuyển} (\en{transition semigroup}).

\begin{congthuc}
\textbf{Tính chất nửa nhóm} (\en{semigroup property}):
\begin{equation}\label{eq:semigroup}
  P(s+t) = P(s)\,P(t) = P(t)\,P(s), \qquad s, t \in \R_+.
\end{equation}
Đây chính là phương trình \term{Chapman--Kolmogorov} phiên bản thời gian liên tục.
\end{congthuc}

% ============================================================================
\section{Ma trận sinh vô cùng bé (Infinitesimal Generator)}
\label{sec:generator}

\textbf{Câu chuyện:} Trong trường hợp rời rạc, ma trận chuyển $P$ cho bạn biết ``xác suất đi đâu sau 1 bước''.
Nhưng trong thời gian liên tục, không có ``1 bước'' --- thời gian chạy liên tục!
Thay vào đó, ta cần biết hệ thống thay đổi \textbf{nhanh thế nào} --- đây là vai trò của ma trận sinh $Q$.

Hãy nghĩ $Q$ như ``tốc độ'' thay vì ``xác suất'':
\begin{itemize}[nosep]
  \item $\lambda_{i,j} > 0$ (ngoài đường chéo): tốc độ chuyển từ $i$ sang $j$. Càng lớn $\Rightarrow$ càng nhanh.
  \item $\lambda_{i,i} < 0$ (đường chéo): tốc độ rời khỏi $i$. $|\lambda_{i,i}|$ lớn $\Rightarrow$ ở $i$ rất ngắn.
\end{itemize}

Bằng cách lấy đạo hàm quan hệ nửa nhóm \eqref{eq:semigroup} theo $t$ rồi cho $t = 0$:
\[
  P'(t) = \lim_{s \to 0} \frac{P(t+s) - P(t)}{s} = P(t) \lim_{s \to 0} \frac{P(s) - P(0)}{s} = P(t)\,Q,
\]
trong đó:

\begin{dinhnghia}[(Ma trận sinh -- \en{Infinitesimal generator})]
\[
  Q := P'(0) = \lim_{s \to 0} \frac{P(s) - P(0)}{s}
\]
được gọi là \term{ma trận sinh vô cùng bé} (\en{infinitesimal generator}) của $(X_t)_{t\in\R_+}$.
\end{dinhnghia}

Ký hiệu $Q = [\lambda_{i,j}]_{i,j\in S}$. Các hàng của $Q$ luôn có tổng bằng 0:
\[
  \sum_{j \in S} \lambda_{i,j} = 0, \qquad i \in S.
\]
Do đó: $\lambda_{i,i} = -\sum_{l \neq i} \lambda_{i,l}$ (phần tử đường chéo luôn âm hoặc bằng 0).

\begin{congthuc}
\textbf{Phương trình Kolmogorov tiến} (\en{forward Kolmogorov equation}):
\begin{equation}\label{eq:forward-kolmogorov}
  P'(t) = P(t)\,Q, \qquad t > 0.
\end{equation}
\textbf{Phương trình Kolmogorov lùi} (\en{backward Kolmogorov equation}):
\begin{equation}\label{eq:backward-kolmogorov}
  P'(t) = Q\,P(t), \qquad t > 0.
\end{equation}
\textbf{Nghiệm} (dùng \term{hàm mũ ma trận} -- \en{matrix exponential}):
\[
  P(t) = P(0)\,e^{tQ} = e^{tQ}, \qquad \text{trong đó } e^{tQ} = I_d + \sum_{n=1}^{\infty} \frac{t^n}{n!} Q^n.
\]
\end{congthuc}

\textbf{Xấp xỉ bậc nhất} khi $h \to 0$:
\[
  P(h) \simeq I_d + hQ + o(h), \qquad h \to 0.
\]

Xác suất chuyển vi mô:
\[
  P_{i,j}(h) = \PP(X_{t+h} = j \mid X_t = i) \simeq
  \begin{cases}
    \lambda_{i,j} h, & i \neq j, \\[3pt]
    1 + \lambda_{i,i} h, & i = j,
  \end{cases}
  \qquad h \to 0.
\]

\textbf{Ý nghĩa vật lý:} Quá trình chuyển từ trạng thái $i$ sang $j \neq i$ giống như quá trình Poisson với cường độ $\lambda_{i,j}$. Thời gian $\tau_{i,j}$ trước khi chuyển từ $i$ sang $j$ có phân phối mũ:
\begin{equation}\label{eq:sojourn-exp}
  \PP(\tau_{i,j} > t) = e^{-\lambda_{i,j} t}, \qquad t \in \R_+, \quad i \neq j.
\end{equation}
và $\E[\tau_{i,j}] = 1/\lambda_{i,j}$.

Thời gian ở trạng thái $i$ (trước khi rời đi bất kỳ đâu):
\[
  \tau_i = \min_{j \in S,\, j \neq i} \tau_{i,j},
\]
là biến mũ với tham số $\sum_{j \neq i} \lambda_{i,j} = -\lambda_{i,i}$, nên $\E[\tau_i] = -1/\lambda_{i,i}$.

\subsection{Ví dụ các ma trận sinh}

\begin{example}[Quá trình Poisson]
Ma trận sinh:
\[
  Q = [\lambda_{i,j}]_{i,j\in\N} = \begin{bmatrix}
    -\lambda & \lambda & 0 & \cdots \\
    0 & -\lambda & \lambda & \cdots \\
    0 & 0 & -\lambda & \cdots \\
    \vdots & \vdots & \vdots & \ddots
  \end{bmatrix}.
\]
Thời gian ở mỗi trạng thái: $\E[\tau_i] = 1/\lambda$ (không phụ thuộc $i$).
\end{example}

\begin{example}[Quá trình sinh-tử trên $\{0,1,\ldots,N\}$]
Ma trận sinh (với $\mu_0 = \lambda_N = 0$):
\[
  Q = \begin{bmatrix}
    -\lambda_0 & \lambda_0 & 0 & \cdots & 0 \\
    \mu_1 & -\lambda_1-\mu_1 & \lambda_1 & \cdots & 0 \\
    \vdots & \ddots & \ddots & \ddots & \vdots \\
    0 & \cdots & \mu_{N-1} & -\lambda_{N-1}-\mu_{N-1} & \lambda_{N-1} \\
    0 & \cdots & 0 & \mu_N & -\mu_N
  \end{bmatrix}.
\]
\end{example}

\begin{example}[Ma trận sinh 5 trạng thái -- tính toán chi tiết]
\label{ex:5state-generator}
Xét quá trình Markov liên tục trên $S = \{0, 1, 2, 3, 4\}$ với ma trận sinh:
\[
  Q = \begin{bmatrix}
    -3 & 2 & 1 & 0 & 0 \\
    1 & -4 & 2 & 1 & 0 \\
    0 & 1 & -3 & 1 & 1 \\
    0 & 0 & 2 & -5 & 3 \\
    0 & 0 & 0 & 4 & -4
  \end{bmatrix}.
\]

\textbf{Đọc thông tin từ $Q$:}
\begin{itemize}[nosep]
  \item Từ trạng thái $0$: chuyển sang $1$ với tốc độ $2$, sang $2$ với tốc độ $1$. Tổng tốc độ rời = $3$.
  \item Thời gian trung bình ở trạng thái $0$: $\E[\tau_0] = 1/3$ (đơn vị thời gian).
  \item Thời gian trung bình ở trạng thái $3$: $\E[\tau_3] = 1/5$ (trạng thái ``bận rộn'' nhất).
  \item Từ trạng thái $4$: chỉ có thể quay lại $3$ (tốc độ $4$), $\E[\tau_4] = 1/4$.
\end{itemize}
\end{example}

\begin{tomtat}
Ma trận sinh $Q$ đóng vai trò tương tự ma trận chuyển $P$ trong trường hợp rời rạc.
Quan hệ: $P(h) \approx I_d + hQ$ cho $h$ nhỏ --- ``rời rạc hóa'' chuỗi liên tục.
Khác biệt chính: phần tử ngoài đường chéo $\lambda_{i,j} \geq 0$ (tốc độ), phần tử đường chéo $\lambda_{i,i} \leq 0$.

\medskip
\textbf{So sánh $P$ (rời rạc) vs $Q$ (liên tục):}
\begin{center}
\renewcommand{\arraystretch}{1.3}
\begin{tabular}{lcc}
\hline
 & \textbf{Ma trận $P$} & \textbf{Ma trận $Q$} \\
\hline
Phần tử ngoài đường chéo & $\geq 0$ (xác suất) & $\geq 0$ (tốc độ) \\
Phần tử đường chéo & $\geq 0$ & $\leq 0$ \\
Tổng mỗi hàng & $= 1$ & $= 0$ \\
Ý nghĩa & ``đi đâu'' & ``đi nhanh thế nào'' \\
\hline
\end{tabular}
\end{center}
\end{tomtat}

% ============================================================================
\section{Chuỗi Markov hai trạng thái liên tục (Two-State Continuous-Time Chain)}
\label{sec:two-state-continuous}

Xét quá trình Markov liên tục trên $S = \{0, 1\}$ với ma trận sinh:
\[
  Q = \begin{bmatrix} -\alpha & \alpha \\ \beta & -\beta \end{bmatrix}, \qquad \alpha, \beta \geq 0.
\]

\textbf{Ví dụ:} Một máy tính có thể ở trạng thái ``hoạt động'' (0) hoặc ``hỏng'' (1).
Máy hỏng với tốc độ $\alpha$ và được sửa với tốc độ $\beta$.

Phương trình Kolmogorov tiến $P'(t) = P(t) Q$ cho hệ phương trình vi phân:
\[
  \begin{cases}
    P'_{0,0}(t) = -\alpha P_{0,0}(t) + \beta P_{0,1}(t), \\
    P'_{0,1}(t) = \alpha P_{0,0}(t) - \beta P_{0,1}(t),
  \end{cases}
\]
với điều kiện đầu $P(0) = I_d$.

\textbf{Nghiệm bằng hàm mũ ma trận:}

Ma trận $Q$ có trị riêng $\lambda_1 = 0$ và $\lambda_2 = -(\alpha + \beta)$.

Chéo hóa và tính $e^{tQ}$:
\begin{equation}\label{eq:Pt-two-state}
  P(t) = e^{tQ} = \frac{1}{\alpha+\beta}
  \begin{bmatrix}
    \beta + \alpha e^{-t(\alpha+\beta)} & \alpha - \alpha e^{-t(\alpha+\beta)} \\
    \beta - \beta e^{-t(\alpha+\beta)} & \alpha + \beta e^{-t(\alpha+\beta)}
  \end{bmatrix}, \quad t > 0.
\end{equation}

\textbf{So sánh với trường hợp rời rạc:} Biểu thức \eqref{eq:Pt-two-state} tương tự công thức rời rạc (Chương~\ref{ch:discrete}), nhưng thay $\lambda_2^n = (1-a-b)^n$ bằng $e^{-t(\alpha+\beta)}$. Sự hội tụ theo hàm mũ thay vì lũy thừa.

\subsection{Phân phối giới hạn}

Khi $t \to \infty$ (với $\alpha > 0$ hoặc $\beta > 0$):
\begin{align}
  \lim_{t\to\infty} \PP(X_t = 1 \mid X_0 = 0) &= \lim_{t\to\infty} \PP(X_t = 1 \mid X_0 = 1) = \frac{\alpha}{\alpha+\beta}, \label{eq:ct-limit-1}\\
  \lim_{t\to\infty} \PP(X_t = 0 \mid X_0 = 0) &= \lim_{t\to\infty} \PP(X_t = 0 \mid X_0 = 1) = \frac{\beta}{\alpha+\beta}. \label{eq:ct-limit-0}
\end{align}

Phân phối giới hạn:
\begin{equation}\label{eq:ct-pi}
  \pi = (\pi_0, \pi_1) = \left(\frac{\beta}{\alpha+\beta},\; \frac{\alpha}{\alpha+\beta}\right).
\end{equation}

\begin{tomtat}
Chuỗi hai trạng thái liên tục có cùng phân phối giới hạn $\pi = (\beta/(\alpha+\beta),\; \alpha/(\alpha+\beta))$ như trường hợp rời rạc (thay $a \to \alpha$, $b \to \beta$).
Sự hội tụ luôn xảy ra trong thời gian liên tục (trừ $\alpha = \beta = 0$), và nhanh hơn khi $\alpha + \beta$ lớn hơn.
\end{tomtat}

% ============================================================================
\section{Phân phối giới hạn và dừng (Limiting and Stationary Distributions)}
\label{sec:ct-stationary}

\textbf{Câu chuyện:} Quay lại ví dụ quán café. Sau nhiều giờ hoạt động, tỷ lệ thời gian quán có $0, 1, 2, \ldots$ khách
sẽ ổn định --- không phụ thuộc vào lúc mở cửa có bao nhiêu khách.
Đây chính là phân phối dừng: ``dân số ổn định'' sau thời gian dài.

\begin{dinhnghia}[(Phân phối dừng liên tục -- \en{Stationary distribution, continuous-time})]
Phân phối $\pi = (\pi_i)_{i\in S}$ được gọi là \term{dừng} (\en{stationary}) nếu $\pi P(t) = \pi$ với mọi $t \in \R_+$.
\end{dinhnghia}

\begin{proposition}[(Mệnh đề 4)]
Phân phối $\pi = (\pi_i)_{i\in S}$ là dừng khi và chỉ khi:
\begin{equation}\label{eq:pi-Q}
  \pi Q = 0.
\end{equation}
\end{proposition}

\textit{Chứng minh.} Nếu $\pi Q = 0$, thì $\pi P(t) = \pi e^{tQ} = \pi \sum_{n=0}^{\infty} \frac{t^n}{n!} Q^n = \pi$.
Ngược lại, lấy đạo hàm $\pi = \pi P(t)$ tại $t = 0$: $0 = \pi P'(0) = \pi Q P(0) = \pi Q$. \qed

\begin{proposition}[(Mệnh đề 5)]
Nếu chuỗi Markov liên tục $(X_t)_{t\in\R_+}$ có phân phối giới hạn
\[
  \pi_j := \lim_{t\to\infty} \PP(X_t = j \mid X_0 = i) = \lim_{t\to\infty} P_{i,j}(t), \qquad j \in S,
\]
không phụ thuộc vào trạng thái ban đầu $i \in S$, thì $\pi Q = 0$, tức $\pi$ là phân phối dừng.
\end{proposition}

\begin{congthuc}
\textbf{Tìm phân phối dừng:} Giải hệ $\pi Q = 0$ kết hợp $\sum_i \pi_i = 1$.

Tương đương: $\sum_{i \in S} \pi_i \lambda_{i,j} = 0$ với mọi $j \in S$.

Lưu ý: điều kiện $\pi Q = 0$ tương đương với $\pi = \pi(I_d + hQ)$ cho mọi $h > 0$, chính là phân phối dừng của chuỗi rời rạc có ma trận chuyển $P(h) \simeq I_d + hQ$.
\end{congthuc}

\subsection{Phân phối dừng của quá trình sinh-tử}

Với quá trình sinh-tử trên $\N$ có ma trận sinh $Q = [\lambda_{i,j}]$, hệ $\pi Q = 0$ cho:
\[
  \begin{cases}
    0 = -\lambda_0 \pi_0 + \mu_1 \pi_1 \\
    0 = \lambda_{j-1}\pi_{j-1} - (\lambda_j + \mu_j)\pi_j + \mu_{j+1}\pi_{j+1}, \quad j \geq 1
  \end{cases}
\]
Giải đệ quy:
\[
  \pi_1 = \frac{\lambda_0}{\mu_1}\pi_0, \qquad
  \pi_{j+1} = \frac{\lambda_j \cdots \lambda_0}{\mu_{j+1} \cdots \mu_1} \pi_0, \qquad j \geq 0,
\]
với $\pi_0$ được xác định bởi $\sum_{i=0}^{\infty} \pi_i = 1$:
\[
  \pi_0 = \frac{1}{\displaystyle\sum_{i=0}^{\infty} \frac{\lambda_{i-1} \cdots \lambda_0}{\mu_i \cdots \mu_1}}.
\]

\textbf{Trường hợp đặc biệt:} Khi $\lambda_i = \lambda$ và $\mu_i = \mu$ với mọi $i \geq 1$:
\[
  \pi_j = \left(1 - \frac{\lambda}{\mu}\right) \left(\frac{\lambda}{\mu}\right)^j, \qquad j \in \N,
\]
với điều kiện $\lambda < \mu$ (phân phối hình học với tham số $\lambda/\mu$).

\begin{example}[Hàng đợi M/M/1]
\label{ex:mm1-queue}
Một quầy phục vụ có khách đến theo Poisson với tốc độ $\lambda$ và thời gian phục vụ mũ với tốc độ $\mu$.
Đây là quá trình sinh-tử với $\lambda_i = \lambda$ và $\mu_i = \mu$ cho mọi $i \geq 1$, $\mu_0 = 0$.

Nếu $\rho = \lambda/\mu < 1$ (hệ thống ổn định):
\[
  \pi_j = (1 - \rho)\rho^j, \qquad j \in \N.
\]

Số khách trung bình trong hệ thống:
\[
  \E[X] = \sum_{j=0}^{\infty} j\, \pi_j = \frac{\rho}{1-\rho} = \frac{\lambda}{\mu - \lambda}.
\]

Ví dụ: nếu $\lambda = 4$ khách/giờ, $\mu = 5$ khách/giờ, thì $\rho = 0.8$ và $\E[X] = 4$ khách.
\end{example}

% ============================================================================
\section{Chuỗi nhúng (The Embedded Chain)}
\label{sec:embedded}

\textbf{Câu chuyện:} Hãy tưởng tượng bạn quay video quán café cả ngày. Chuỗi Markov liên tục là toàn bộ video.
Nhưng nếu bạn chỉ chụp ảnh mỗi khi có sự thay đổi (khách đến hoặc đi),
bạn thu được một ``bộ ảnh'' --- đây chính là \term{chuỗi nhúng}.

Chuỗi nhúng ``bỏ qua thời gian'', chỉ ghi lại \emph{trạng thái nào} xuất hiện, không quan tâm \emph{bao lâu} ở mỗi trạng thái.

\begin{dinhnghia}[(Chuỗi nhúng -- \en{Embedded chain})]
Xét dãy $(T_n)_{n\in\N}$ các thời điểm nhảy của quá trình Markov liên tục $(X_t)_{t\in\R_+}$:
\[
  T_1 = \inf\{t > 0 : X_t \neq X_0\}, \qquad
  T_{n+1} = \inf\{t > T_n : X_t \neq X_{T_n}\}, \quad n \in \N.
\]
\term{Chuỗi Markov nhúng} (\en{embedded Markov chain}) là chuỗi rời rạc $(Z_n)_{n\in\N}$ với:
\[
  Z_0 = X_0, \qquad Z_n = X_{T_n}, \quad n \in \N.
\]
\end{dinhnghia}

\subsection{Ma trận chuyển của chuỗi nhúng}

Với quá trình sinh-tử, ma trận chuyển của chuỗi nhúng:
\[
  P_{i,i+1} = \frac{\lambda_i}{\lambda_i + \mu_i}, \qquad
  P_{i,i-1} = \frac{\mu_i}{\lambda_i + \mu_i}, \qquad i \in \N.
\]

Điều này là do:
\[
  \PP(\tau_{i,i+1} < \tau_{i,i-1}) = \frac{\lambda_i}{\lambda_i + \mu_i},
\]
theo tính chất cạnh tranh giữa hai biến mũ độc lập.

\textbf{Tổng quát:} Với chuỗi Markov liên tục có ma trận sinh $Q = [\lambda_{i,j}]$,
ma trận chuyển của chuỗi nhúng là:
\begin{equation}\label{eq:embedded-general}
  \hat{P}_{i,j} = \frac{\lambda_{i,j}}{-\lambda_{i,i}} = \frac{\lambda_{i,j}}{\sum_{l \neq i} \lambda_{i,l}}, \qquad i \neq j.
\end{equation}

\begin{example}[Từ ma trận $Q$ đến chuỗi nhúng -- 5 trạng thái]
Xét quá trình ở Ví dụ~\ref{ex:5state-generator} với ma trận sinh:
\[
  Q = \begin{bmatrix}
    -3 & 2 & 1 & 0 & 0 \\
    1 & -4 & 2 & 1 & 0 \\
    0 & 1 & -3 & 1 & 1 \\
    0 & 0 & 2 & -5 & 3 \\
    0 & 0 & 0 & 4 & -4
  \end{bmatrix}.
\]

Ma trận chuyển của chuỗi nhúng $\hat{P}$:
\[
  \hat{P} = \begin{bmatrix}
    0 & 2/3 & 1/3 & 0 & 0 \\
    1/4 & 0 & 2/4 & 1/4 & 0 \\
    0 & 1/3 & 0 & 1/3 & 1/3 \\
    0 & 0 & 2/5 & 0 & 3/5 \\
    0 & 0 & 0 & 1 & 0
  \end{bmatrix}.
\]

\textbf{Cách tính:} Hàng $i$ của $\hat{P}$ = hàng $i$ của $Q$ (bỏ phần tử đường chéo), chia cho $|\lambda_{i,i}|$.

Ví dụ hàng 1 (trạng thái 1): $\lambda_{1,0} = 1$, $\lambda_{1,2} = 2$, $\lambda_{1,3} = 1$, $|\lambda_{1,1}| = 4$.
\[
  \hat{P}_{1,0} = \frac{1}{4}, \quad \hat{P}_{1,2} = \frac{2}{4}, \quad \hat{P}_{1,3} = \frac{1}{4}.
\]
\end{example}

\begin{tomtat}
Chuỗi nhúng ``rời rạc hóa'' quá trình liên tục bằng cách chỉ ghi lại các trạng thái tại thời điểm nhảy.
Ma trận chuyển của chuỗi nhúng chỉ phụ thuộc vào \emph{tỷ lệ} các tốc độ, không phụ thuộc vào tốc độ tuyệt đối.

Để phục hồi quá trình liên tục từ chuỗi nhúng, cần thêm thông tin về \emph{thời gian ở mỗi trạng thái}
$\tau_i \sim \text{Exp}(-\lambda_{i,i})$.
\end{tomtat}

% ============================================================================
\section{Xác suất hấp thụ (Absorption Probabilities)}
\label{sec:absorption-prob}

\textbf{Câu chuyện:} Hãy tưởng tượng một con kiến bắt đầu ở vị trí $i$ trên một thanh gỗ.
Hai đầu thanh gỗ ($0$ và $N$) là ``tường'' --- khi chạm tường, kiến dừng lại vĩnh viễn.
Câu hỏi: ``Con kiến sẽ chạm tường nào? Với xác suất bao nhiêu?''

Đây là bài toán \term{xác suất hấp thụ} cho quá trình sinh-tử liên tục.
Vì chuỗi nhúng quyết định ``đi đâu'' (không quan tâm ``bao lâu''),
xác suất hấp thụ của quá trình liên tục \emph{bằng} xác suất hấp thụ của chuỗi nhúng.

\subsection{Thiết lập bài toán}

Xét quá trình sinh-tử trên $S = \{0, 1, \ldots, N\}$ với:
\begin{itemize}[nosep]
  \item Trạng thái $0$ và $N$ là hấp thụ: $\lambda_0 = 0$ (không sinh khi ở 0) và $\mu_N = 0$ (không tử khi ở $N$).
  \item Khi ở trạng thái $i$ ($1 \leq i \leq N-1$): sinh với tốc độ $\lambda_i$, tử với tốc độ $\mu_i$.
\end{itemize}

Đặt $g_0(i) = \PP(\text{bị hấp thụ vào trạng thái } 0 \mid X_0 = i)$.

\subsection{Phương trình cho xác suất hấp thụ}

Vì xác suất hấp thụ chỉ phụ thuộc chuỗi nhúng, ta áp dụng phân tích bước đầu tiên trên chuỗi nhúng:
\begin{equation}\label{eq:absorption-fsa}
  g_0(i) = \hat{P}_{i,i-1}\, g_0(i-1) + \hat{P}_{i,i+1}\, g_0(i+1),
  \qquad 1 \leq i \leq N-1,
\end{equation}
với điều kiện biên $g_0(0) = 1$ và $g_0(N) = 0$.

Thay $\hat{P}_{i,i-1} = \frac{\mu_i}{\lambda_i + \mu_i}$ và $\hat{P}_{i,i+1} = \frac{\lambda_i}{\lambda_i + \mu_i}$:
\begin{equation}\label{eq:absorption-bd}
  g_0(i) = \frac{\mu_i}{\lambda_i + \mu_i}\, g_0(i-1) + \frac{\lambda_i}{\lambda_i + \mu_i}\, g_0(i+1).
\end{equation}

Đây là phương trình ``con bạc'' (\en{gambler's equation}) ở dạng liên tục.

\subsection{Nghiệm tường minh}

Đặt $\rho_i = \mu_i / \lambda_i$ (tỷ số tử/sinh). Viết lại \eqref{eq:absorption-bd}:
\[
  g_0(i) - g_0(i-1) = \frac{\lambda_i}{\mu_i}\left[g_0(i+1) - g_0(i)\right],
\]
hay:
\[
  g_0(i+1) - g_0(i) = \rho_i \left[g_0(i) - g_0(i-1)\right].
\]

Đặt $d_i = g_0(i) - g_0(i-1)$. Ta có $d_{i+1} = \rho_i\, d_i$, suy ra:
\[
  d_k = d_1 \prod_{j=1}^{k-1} \rho_j = d_1 \prod_{j=1}^{k-1} \frac{\mu_j}{\lambda_j}, \qquad k \geq 2.
\]

Từ $g_0(0) = 1$ và $g_0(N) = 0$:
\[
  g_0(N) - g_0(0) = \sum_{k=1}^{N} d_k = -1.
\]

Đặt $R_k = \prod_{j=1}^{k-1} \frac{\mu_j}{\lambda_j}$ (với $R_1 = 1$). Ta được $d_1 \sum_{k=1}^{N} R_k = -1$, nên:

\begin{congthuc}
\textbf{Xác suất hấp thụ vào trạng thái $0$:}
\begin{equation}\label{eq:absorption-formula}
  g_0(i) = \frac{\displaystyle\sum_{k=i+1}^{N} R_k}{\displaystyle\sum_{k=1}^{N} R_k},
  \qquad \text{với } R_k = \prod_{j=1}^{k-1} \frac{\mu_j}{\lambda_j},\; R_1 = 1.
\end{equation}
\end{congthuc}

\textbf{Trường hợp đặc biệt} $\lambda_i = \lambda$, $\mu_i = \mu$ (tốc độ không đổi):

Nếu $\rho = \mu/\lambda \neq 1$: $R_k = \rho^{k-1}$, suy ra:
\begin{equation}\label{eq:absorption-constant}
  g_0(i) = \frac{\rho^i - \rho^N}{1 - \rho^N}.
\end{equation}

Nếu $\rho = 1$ (tức $\lambda = \mu$):
\[
  g_0(i) = 1 - \frac{i}{N}.
\]

\begin{example}[Xác suất hấp thụ với $N = 4$]
\label{ex:absorption-N4}
Xét quá trình sinh-tử trên $\{0, 1, 2, 3, 4\}$ với $\lambda_i = 2$, $\mu_i = 3$ cho $1 \leq i \leq 3$.

Ta có $\rho = \mu/\lambda = 3/2$. Áp dụng \eqref{eq:absorption-constant}:
\[
  g_0(i) = \frac{(3/2)^i - (3/2)^4}{1 - (3/2)^4}
  = \frac{(3/2)^i - 81/16}{1 - 81/16}
  = \frac{(3/2)^i - 81/16}{-65/16}.
\]

Tính cụ thể:
\begin{align*}
  g_0(1) &= \frac{3/2 - 81/16}{-65/16} = \frac{24/16 - 81/16}{-65/16} = \frac{-57/16}{-65/16} = \frac{57}{65} \approx 0.877, \\
  g_0(2) &= \frac{9/4 - 81/16}{-65/16} = \frac{36/16 - 81/16}{-65/16} = \frac{45}{65} = \frac{9}{13} \approx 0.692, \\
  g_0(3) &= \frac{27/8 - 81/16}{-65/16} = \frac{54/16 - 81/16}{-65/16} = \frac{27}{65} \approx 0.415.
\end{align*}

\textbf{Nhận xét:} Vì $\mu > \lambda$ (tốc độ tử $>$ tốc độ sinh), quá trình ``kéo về 0'' mạnh hơn,
nên xác suất chạm 0 khá cao ngay cả từ trạng thái 3.
\end{example}

% ============================================================================
\section{Thời gian hấp thụ trung bình (Mean Absorption Times)}
\label{sec:absorption-time}

\textbf{Câu chuyện:} Ta đã biết con kiến sẽ chạm tường. Nhưng \textbf{mất bao lâu}?
Trong thời gian liên tục, đây là câu hỏi quan trọng vì thời gian có ý nghĩa thực tế
(ví dụ: ``Trung bình bao lâu thì hệ thống tắt?'').

Khác với chuỗi rời rạc (đếm số bước), ở đây ta cần cộng các \emph{thời gian lưu trú} (\en{sojourn times})
tại mỗi trạng thái.

\subsection{Thiết lập}

Đặt $\eta(i) = \E[T_{\text{abs}} \mid X_0 = i]$, trong đó $T_{\text{abs}}$ là thời gian đến khi bị hấp thụ
(chạm $0$ hoặc $N$).

Với quá trình sinh-tử, phân tích bước đầu tiên \emph{trên chuỗi nhúng} cho:
\begin{equation}\label{eq:mean-abs-embedded}
  \eta(i) = \E[\tau_i] + \hat{P}_{i,i-1}\, \eta(i-1) + \hat{P}_{i,i+1}\, \eta(i+1),
  \qquad 1 \leq i \leq N-1,
\end{equation}
với điều kiện biên $\eta(0) = \eta(N) = 0$.

Ở đây $\E[\tau_i] = \frac{1}{\lambda_i + \mu_i}$ là thời gian trung bình ở trạng thái $i$ trước khi nhảy.

\textbf{Sự khác biệt so với rời rạc:} Trong trường hợp rời rạc, ta có ``$+1$'' (mỗi bước mất 1 đơn vị thời gian).
Ở đây, ta thay bằng ``$+\E[\tau_i]$'' --- mỗi bước mất thời gian $\text{Exp}(\lambda_i + \mu_i)$.

\subsection{Nghiệm}

Thay $\hat{P}_{i,i\pm 1}$ và $\E[\tau_i]$ vào \eqref{eq:mean-abs-embedded}:
\[
  \eta(i) = \frac{1}{\lambda_i + \mu_i} + \frac{\mu_i}{\lambda_i + \mu_i}\, \eta(i-1) + \frac{\lambda_i}{\lambda_i + \mu_i}\, \eta(i+1).
\]

Nhân cả hai vế với $(\lambda_i + \mu_i)$:
\[
  (\lambda_i + \mu_i)\eta(i) = 1 + \mu_i\, \eta(i-1) + \lambda_i\, \eta(i+1).
\]

Viết lại:
\begin{equation}\label{eq:mean-abs-recurrence}
  \lambda_i[\eta(i) - \eta(i+1)] = 1 + \mu_i[\eta(i-1) - \eta(i)].
\end{equation}

Đặt $\Delta_i = \eta(i) - \eta(i+1)$. Ta có quan hệ đệ quy:
\[
  \lambda_i \Delta_i = 1 - \mu_i \Delta_{i-1},
\]
hay:
\[
  \Delta_i = \frac{1}{\lambda_i} - \frac{\mu_i}{\lambda_i} \Delta_{i-1}
  = \frac{1}{\lambda_i} + \rho_i \Delta_{i-1}.
\]

\begin{example}[Thời gian hấp thụ trung bình với $N = 4$]
Tiếp tục Ví dụ~\ref{ex:absorption-N4} với $\lambda_i = 2$, $\mu_i = 3$.

Thời gian trung bình ở mỗi trạng thái: $\E[\tau_i] = 1/(\lambda_i + \mu_i) = 1/5$ giờ = 12 phút.

Hệ phương trình (nhân cả hai vế với 5):
\[
  \begin{cases}
    5\eta(1) = 1 + 3\eta(0) + 2\eta(2) = 1 + 2\eta(2), \\
    5\eta(2) = 1 + 3\eta(1) + 2\eta(3), \\
    5\eta(3) = 1 + 3\eta(2) + 2\eta(4) = 1 + 3\eta(2).
  \end{cases}
\]

Từ PT3: $\eta(3) = \frac{1 + 3\eta(2)}{5}$.

Thay vào PT2: $5\eta(2) = 1 + 3\eta(1) + 2 \cdot \frac{1 + 3\eta(2)}{5}$
$\Rightarrow 25\eta(2) = 5 + 15\eta(1) + 2 + 6\eta(2)$
$\Rightarrow 19\eta(2) = 7 + 15\eta(1)$.

Từ PT1: $\eta(1) = \frac{1 + 2\eta(2)}{5}$.

Thay: $19\eta(2) = 7 + 15 \cdot \frac{1 + 2\eta(2)}{5} = 7 + 3(1 + 2\eta(2)) = 10 + 6\eta(2)$.

Suy ra: $13\eta(2) = 10$, $\eta(2) = 10/13$.

Khi đó: $\eta(1) = \frac{1 + 20/13}{5} = \frac{33/13}{5} = \frac{33}{65}$,
$\eta(3) = \frac{1 + 30/13}{5} = \frac{43/13}{5} = \frac{43}{65}$.

\textbf{Kết quả:}
\[
  \eta(1) = \frac{33}{65} \approx 0.508, \quad
  \eta(2) = \frac{10}{13} \approx 0.769, \quad
  \eta(3) = \frac{43}{65} \approx 0.662.
\]

(Đơn vị: giờ, nếu $\lambda, \mu$ tính theo giờ.)

\textbf{Nhận xét:} $\eta(2) > \eta(3) > \eta(1)$. Trạng thái $2$ (ở giữa) mất lâu nhất
vì nó xa cả hai biên hấp thụ.
\end{example}

% ============================================================================
\section{Quan hệ giữa phân phối dừng liên tục và rời rạc}
\label{sec:continuous-vs-discrete}

Một câu hỏi tự nhiên: phân phối dừng $\pi$ của quá trình liên tục
và phân phối dừng $\hat{\pi}$ của chuỗi nhúng có quan hệ gì?

\begin{dinhly}[(Quan hệ $\pi$ và $\hat{\pi}$)]
Nếu $\pi$ là phân phối dừng của quá trình Markov liên tục (tức $\pi Q = 0$),
và $\hat{\pi}$ là phân phối dừng của chuỗi nhúng (tức $\hat{\pi} \hat{P} = \hat{\pi}$), thì:
\begin{equation}\label{eq:pi-hat-relation}
  \pi_i \propto \frac{\hat{\pi}_i}{-\lambda_{i,i}}, \qquad i \in S.
\end{equation}
Tức là: $\pi_i = \frac{\hat{\pi}_i / (-\lambda_{i,i})}{\sum_j \hat{\pi}_j / (-\lambda_{j,j})}$.
\end{dinhly}

\textbf{Ý nghĩa trực quan:} Phân phối dừng liên tục $\pi_i$ tỷ lệ với cả hai yếu tố:
\begin{enumerate}[nosep]
  \item Tần suất ``ghé thăm'' trạng thái $i$ (đo bằng $\hat{\pi}_i$).
  \item Thời gian ở lại trạng thái $i$ mỗi lần ghé (đo bằng $1/(-\lambda_{i,i})$).
\end{enumerate}
Trạng thái được ghé thăm thường xuyên VÀ giữ chân lâu sẽ có $\pi_i$ lớn.

% ============================================================================
\section{Bài tập (Exercises)}
\label{sec:ch9-exercises}

\begin{exercise}
Xét quá trình sinh-tử trên $\{0, 1, 2, 3\}$ với tốc độ sinh $\lambda_i = 2$ và tốc độ tử $\mu_i = 3$ cho $1 \leq i \leq 2$, $\lambda_0 = 2$, $\mu_3 = 3$, $\lambda_3 = \mu_0 = 0$.
\begin{enumerate}[nosep]
  \item Viết ma trận sinh $Q$.
  \item Tìm phân phối dừng $\pi$.
  \item Viết ma trận chuyển của chuỗi nhúng.
\end{enumerate}
\end{exercise}

\begin{exercise}
Cho chuỗi Markov liên tục hai trạng thái với $\alpha = 3$ và $\beta = 6$.
\begin{enumerate}[nosep]
  \item Tính $P(t)$ cho mọi $t > 0$.
  \item Tính thời gian trung bình ở mỗi trạng thái.
  \item Tìm phân phối giới hạn.
\end{enumerate}
\end{exercise}

\begin{exercise}[Hàng đợi M/M/1]
Khách đến quầy theo Poisson với tốc độ $\lambda = 3$ khách/giờ.
Thời gian phục vụ mỗi khách là biến mũ với tốc độ $\mu = 5$ phục vụ/giờ.
\begin{enumerate}[nosep]
  \item Viết ma trận sinh $Q$ (vô hạn chiều).
  \item Tìm phân phối dừng. Hệ thống có ổn định không?
  \item Tính số khách trung bình trong hệ thống.
  \item Tính xác suất hệ thống rỗng (không có khách nào).
\end{enumerate}
\end{exercise}

\begin{exercise}[Quá trình điện tín -- \en{Telegraph process}]
Tín hiệu điện tín chuyển giữa ``bật'' và ``tắt'' theo chuỗi Markov liên tục hai trạng thái
với tốc độ bật$\to$tắt là $\alpha = 2$ và tắt$\to$bật là $\beta = 3$.
\begin{enumerate}[nosep]
  \item Tìm xác suất tín hiệu đang ``bật'' tại thời điểm $t$, biết ban đầu ``bật''.
  \item Tìm phân phối giới hạn. Bao nhiêu phần trăm thời gian tín hiệu ``bật''?
  \item Thời gian trung bình một chu kỳ ``bật'' kéo dài bao lâu?
\end{enumerate}
\end{exercise}

\begin{exercise}[Xác suất hấp thụ]
Xét quá trình sinh-tử trên $\{0, 1, 2, 3, 4, 5\}$ với $\lambda_i = 1$ và $\mu_i = 2$ cho $1 \leq i \leq 4$,
trạng thái $0$ và $5$ hấp thụ.
\begin{enumerate}[nosep]
  \item Tính xác suất bị hấp thụ vào trạng thái $0$ xuất phát từ $i = 1, 2, 3, 4$.
  \item So sánh với trường hợp $\lambda_i = \mu_i = 1$ (tốc độ bằng nhau).
\end{enumerate}
\end{exercise}

\begin{exercise}[Quá trình Yule -- \en{Yule process}]
Quá trình Yule là quá trình sinh thuần túy với $\lambda_i = i\lambda$ (tốc độ sinh tỷ lệ với số cá thể hiện tại).
Bắt đầu với $X_0 = 1$.
\begin{enumerate}[nosep]
  \item Tính thời gian trung bình ở trạng thái $i$ (tức $\E[\tau_i]$).
  \item Tính thời gian trung bình để quần thể tăng từ $1$ lên $n$ cá thể.
  \item Giải thích tại sao quá trình Yule mô hình hóa sự tăng trưởng dân số.
\end{enumerate}
\end{exercise}

\begin{exercise}[Thời gian hấp thụ trung bình]
Xét quá trình sinh-tử trên $\{0, 1, 2, 3\}$ với $\lambda_i = 1$, $\mu_i = 1$ cho $i = 1, 2$,
trạng thái $0$ và $3$ hấp thụ.
\begin{enumerate}[nosep]
  \item Tính thời gian hấp thụ trung bình $\eta(1)$ và $\eta(2)$.
  \item So sánh với thời gian hấp thụ trong chuỗi rời rạc tương ứng (gambler's ruin với $p = q = 1/2$, $N = 3$).
\end{enumerate}
\end{exercise}
